\documentclass[11pt]{article}
\usepackage[margin=1in]{geometry}
\usepackage{amsmath,amssymb,amsthm,hyperref}
\newtheorem{theorem}{Theorem}
\newtheorem{lemma}{Lemma}
\newcommand{\Q}{\mathbb Q}
\newcommand{\Z}{\mathbb Z}
\newcommand{\R}{\mathbb R}
\title{A rational counterexample to the nonreducibility characterization}
\author{Principia Math In Progress}
\date{27 September 2026}
\begin{document}
\maketitle

\paragraph{Task 1: the source question.}
Conjecture 5.3 of Kiss and Somlai \cite{KS} concerns associative,
idempotent, strictly coordinatewise monotone operations on a totally
ordered Abelian group $(X,+)$. It proposes that a nonreducible operation
must have odd arity and be conjugate, by a monotone bijection $X\to X$,
to the alternating sum formed using the given group addition.
Here reducibility means iteration of an associative binary operation on
the same set; no monotonicity of that binary operation is assumed.
The printed sum starts with a minus sign. We refute the substantive
interpretation starting with a plus sign, as in the source's Example 4.10.
Thus the counterexample does not depend on that typographical issue.

\begin{theorem}
There is an associative, idempotent operation $F:\Q^3\to\Q$, strictly
increasing in its first and third variables and strictly decreasing in its
second variable, which is nonreducible and for which no bijection
$g:\Q\to\Q$ satisfies
\begin{equation}\label{eq:claimed}
 F(x,y,z)=g^{-1}\bigl(g(x)-g(y)+g(z)\bigr).
\end{equation}
In particular, the characterization in Conjecture 5.3 fails for the
usual ordered Abelian group $(\Q,+)$, even after correction of the sign.
\end{theorem}

\begin{lemma}\label{lem:order}
Let $H=\Z+\sqrt2\Z\subset\R$ with its usual order. There exists an
increasing bijection $h:\Q\to H$ with $h(0)=0$ and $h(1)=1$.
However, $1/2\notin H$.
\end{lemma}
\begin{proof}
The set $H$ is a countable additive subgroup, unbounded in both
directions. To see that it is dense, divide $[0,1)$ into $N$ equal
half-open intervals and consider the $N+1$ fractional parts of
$j\sqrt2$, $0\le j\le N$. Two lie in the same interval; their nonzero
difference, with sign chosen positive, gives $d\in H$ with $0<d<1/N$.
For any real $a<b$, choose such a $d<b-a$. Then
$(\lfloor a/d\rfloor+1)d$ lies in $(a,b)\cap H$.

For completeness, construct $h$ by back-and-forth. Enumerate both
countable sets and start with $0\mapsto0$, $1\mapsto1$. At alternate
steps, take the first rational not yet in the finite domain and map it
to an unused element of $H$ between the images of its nearest domain
neighbors; when a neighbor is missing use the corresponding unbounded
interval. At the other steps, take the first element of $H$ not yet in
the range and choose an unused rational in the corresponding interval.
Density and absence of endpoints make each choice possible. The union
is increasing, defined on every rational, and onto $H$.

Finally, $a+b\sqrt2=1/2$ with $a,b\in\Z$ would force $b=0$ by
irrationality of $\sqrt2$, and then $a=1/2$, a contradiction.
\end{proof}

\begin{proof}[Proof of the theorem]
Fix $h$ from Lemma~\ref{lem:order} and define
\begin{equation}\label{eq:construction}
 F(x,y,z)=h^{-1}\bigl(h(x)-h(y)+h(z)\bigr).
\end{equation}
Closure of $H$ under addition and subtraction makes this a total
operation on $\Q$. Its strict monotonicity in each coordinate and its
idempotency follow immediately from the increasing bijection $h$.
Writing $T(u,v,w)=u-v+w$ on $H$, all three ternary associativity
expressions
\[
 T(T(u_1,u_2,u_3),u_4,u_5),\quad
 T(u_1,T(u_2,u_3,u_4),u_5),\quad
 T(u_1,u_2,T(u_3,u_4,u_5))
\]
equal $u_1-u_2+u_3-u_4+u_5$. Transport through $h$ proves
associativity of $F$ in the source's precise sense.

Suppose an associative binary operation on $\Q$ reduced $F$.
Transporting it through $h$ would give an associative operation $*$ on
$H$ such that
\[
 (u*v)*w=u*(v*w)=u-v+w.
\]
Let $e=0*0$. For every $z\in H$,
\[
 e*z=(0*0)*z=T(0,0,z)=z,\qquad
 z*e=z*(0*0)=T(z,0,0)=z.
\]
Thus $e$ is a two-sided identity. Consequently
$T(e,x,e)=(e*x)*e=x$ for all $x\in H$. But the definition of $T$
gives $T(e,x,e)=2e-x$. Hence $2(x-e)=0$ for every $x\in H$,
contradicting, for instance, $x=e+1$. This proves nonreducibility
without any assumptions on a proposed binary reduction except
associativity.

On the other hand, equation $F(x,0,x)=1$ has no solution: by
\eqref{eq:construction} and the normalization of $h$, it is equivalent
to $2h(x)=1$, impossible by Lemma~\ref{lem:order}.
If a bijection $g:\Q\to\Q$ satisfied \eqref{eq:claimed}, the rational
\[
 x=g^{-1}\!\left(\frac{g(0)+g(1)}2\right)
\]
would solve precisely that equation. This contradiction excludes every
bijection $g$, and therefore every monotone one.
\end{proof}

\paragraph{Scope.}
The group addition in the conjecture is fixed. Formula
\eqref{eq:construction} transports addition from the different ordered
group $H$, which has the same order type as $\Q$ but different
divisibility properties. The theorem disproves the universally
quantified characterization at $n=3$. It makes no claim about the
restricted case $X=\R$, or about a formulation allowing a different
target group for the conjugating map. The original printed minus-first
sum cannot represent an idempotent operation on $\Q$ either: on the
diagonal it would require $-g(x)=g(x)$ for every $x$.

\paragraph{Proposed follow-up tasks (not claimed solved here).}
\begin{enumerate}
\setcounter{enumi}{1}
\item Extend the counterexample to every odd arity $n\ge3$, proving
all associativity identities and nonreducibility without invoking
the source's reducibility theorems.
\item Classify when the ternary operations $u-v+w$ on two dense
additive subgroups of $\R$ are conjugate by an increasing bijection;
determine whether a translation followed by positive scaling is
necessary and sufficient.
\item Resolve the corrected Conjecture 5.3 for $X=\R$ under strict
coordinatewise monotonicity alone, including the proposed automatic
continuity conclusion.
\item Construct infinitely many pairwise nonconjugate associative,
idempotent, strictly monotone ternary operations on $\Q$, distinguishing
them by explicit equations expressing prime divisibility.
\end{enumerate}

\paragraph{Verification.}
The companion exact verifier checks the universal linear identities,
the binary-reduction obstruction, and the midpoint calculation.
Existence of the infinite bijection and its order properties are proved
above, not inferred from finite experiments. Review is by the authoring
agent; this is not independent peer review or a proof-assistant certificate.

\begin{thebibliography}{1}
\bibitem{KS} G. Kiss and G. Somlai,
\emph{Associative idempotent nondecreasing functions are reducible},
Semigroup Forum \textbf{98} (2019), 140--153.
\href{https://doi.org/10.1007/s00233-018-9973-y}{doi:10.1007/s00233-018-9973-y}.
Conjecture 5.3 and Example 4.10; also
\href{https://arxiv.org/abs/1707.04341v3}{arXiv:1707.04341v3}.
The 16-page postprint deposited at
\href{https://orbilu.uni.lu/handle/10993/34510}{ORBilu 10993/34510}
retains the conjecture on page 15.
\end{thebibliography}
\end{document}
